\section{粮食与油料：玉米、大豆及压榨链}
\label{sec:grain}

如果说猪周期是农业需求侧的节拍器，那么\hl{粮食与油料就是被这个节拍器驱动、但有自己的旋律}的
另一条主线。中国农业的“粮食—饲料—养殖”链条中，玉米与大豆是两个枢纽品种：前者是国内唯一
能够横跨饲料与深加工两大需求的大宗粮食品种，后者则以超过 1 亿吨的年进口量把中国农业与
南美、北美的天气和贸易政策直接绑定。

\subsection{产业链结构：两类枢纽品种}
\label{sec:grain-structure}

\begin{center}
\begin{tikzpicture}[font=\footnotesize,
  n/.style={rectangle, rounded corners=2pt, draw=Grain!80!black, line width=0.8pt,
            fill=Grain!12, align=center, inner sep=3.5pt, minimum height=1.05cm},
  e/.style={-{Stealth[length=3.5pt]}, line width=0.9pt, InkGray},
  lab/.style={font=\scriptsize, InkGray, align=center}]
  % 玉米链
  \node[n, text width=1.5cm] (cg) at (0,2.6) {玉米\\30123 万吨};
  \node[n, text width=1.9cm] (cf) at (2.9,2.6) {饲料用\\65\%$\sim$70\%};
  \node[n, text width=1.7cm] (ci) at (5.9,2.6) {深加工\\25\%$\sim$30\%};
  \node[n, text width=1.7cm] (cl) at (8.9,2.6) {生猪养殖\\饲用中约 50\%};
  \draw[e] (cg) -- (cf);
  \draw[e] (cg) -- (ci);
  \draw[e] (cf) -- (cl);
  % 大豆链
  \node[n, text width=1.9cm] (sg) at (0,0.4) {大豆\\进口 11183 + 国产 2091};
  \node[n, text width=1.6cm] (sp) at (3.1,0.4) {压榨\\10150 万吨};
  \node[n, text width=1.5cm] (sm) at (6.2,0.4) {豆粕\\约 80\%};
  \node[n, text width=1.5cm] (so) at (8.6,0.4) {豆油\\约 18\%};
  \node[n, text width=1.7cm] (sj) at (6.2,-1.8) {生猪、家禽\\饲料（刚性需求）};
  \draw[e] (sg) -- (sp);
  \draw[e] (sp) -- (sm);
  \draw[e] (sp) -- (so);
  \draw[e] (sm) -- (sj);
  \node[lab, anchor=west] at (-1.6,3.5) {（a）玉米：两条需求腿};
  \node[lab, anchor=west] at (-1.6,1.3) {（b）大豆：压榨链};
\end{tikzpicture}
\end{center}

\textbf{玉米的两条需求腿。}玉米饲用消费占总消费量的 \hl{65\%$\sim$70\%}，其中生猪养殖消耗
约占饲用消费的 50\%、禽料约 40\%；深加工年耗玉米约 7800 万吨、占 25\%$\sim$30\%，
全国深加工产能已超过 1.25 亿吨\cite{lswz2025corn}。这一结构决定了玉米价格同时受
\emph{养殖产能周期}与\emph{深加工开工率}影响，而两者的周期长度完全不同：
养殖产能是 3$\sim$4 年的长周期，深加工利润则是 1 年以内的短周期。

\textbf{大豆的压榨链。}大豆压榨产出约 80\% 的豆粕与 18\% 的豆油，豆粕几乎全部用于饲料。
2025/26 年度中国大豆压榨量约 10150 万吨\cite{casde123}，压榨产能 18963 万吨、近五年
复合增长率 2.33\%\cite{zhuochuang2026soy}。由于豆粕需求来自生猪与家禽存栏，
\hl{压榨量与猪周期同向}，但豆粕的\emph{价格}由进口成本（CBOT + 升贴水 + 汇率）决定，
这是理解“量与价背离”的关键。

\subsection{供给端：产量、进口与政策}
\label{sec:grain-supply}

\subsubsection{玉米：连续丰产与进口政策转向}

\begin{table}[htbp]
\centering
\small
\begin{tabular}{@{}llr@{}}\toprule
指标 & 最新值 & 备注 \\
\midrule
播种面积 & 44960.8 千公顷 & 2025 年（玉米） \\
产量 & 30123.5 万吨 & 2025 年，单产 6700.0 公斤/公顷 \\
进口量 & 264.77 万吨 & 2025 年日历年，同比 -80.8\% \\
进口量 & 182.4 万吨 & 2024/25 年度（2024 年 10 月---2025 年 9 月），同比 -92\% \\
进口配额 & 720 万吨 & 2004 年至今未变，配额内关税 1\%、配额外 65\% \\
替代品进口 & 2301 万吨 & 2024/25 年度高粱、大麦、DDGS、木薯干合计，同比 -54\% \\
\bottomrule
\end{tabular}
\caption{中国玉米供给端关键指标（2024---2025）}
\label{tab:corn-supply}
\end{table}

玉米供给端最重要的事件是\hl{进口的断崖式退出}：在 2021---2023 年连续三年进口超过
2000 万吨之后，2024 年进口回落至 1364 万吨（-49.7\%），
2024/25 年度仅 182.4 万吨（-92\%，近十四年最低），其中美玉米仅 3.9 万吨（-98.7\%）、
巴西玉米 72.5 万吨（-95.2\%）\cite{xinhu2025corn}。其直接原因是
2024 年 4 月前后进口政策收紧，月度进口量从百万吨级降至 20 万吨以下\cite{xinhu2025corn}。

这一转变的政策含义是：\hl{国内玉米价格的政策底由“进口成本”变成了“国内种植成本 + 补贴”}。
进口配额 20 余年稳定在 720 万吨、配额外 65\% 的关税，意味着只要进口量被行政性压低，
国内玉米价格就失去了外部的天花板与地板，\emph{价格更多由国内产量与库存决定}。
2026/27 年度预测产量 \num{30600} 万吨（+1.6\%）\cite{casde119}；
2025/26 年度饲用消费估计值上调至 \num{21090} 万吨、工业消费下调至 \num{8100} 万吨\cite{casde119}，
供需处于紧平衡状态。

\subsubsection{大豆：创纪录进口与来源结构重构}

2025 年中国大豆进口 \hl{11183 万吨}（同比 +6.47\%，历史新高），
其中巴西 8232 万吨（占 73.61\%，同比提升 2.54 个百分点）、
美国 1682 万吨（占 15.04\%，同比下滑 6.03 个百分点，为 2015 年有数据以来最低）、
阿根廷 789 万吨（+92.4\%）；进口均价 450 美元/吨、同比下降 10.35\%\cite{zhuochuang2026soy,smm2026soy}。
2025 年 12 月自美国进口大豆为零，为连续第四个月零进口\cite{smm2026soy}。
国产大豆 2025 年产量 2091 万吨（+1.3\%）\cite{casde119}，仅相当于进口量的 19\%，
\hl{大豆是农业领域进口依赖度最高、也最容易被贸易政策扰动的品种}。

\begin{keybox}[供给端的三种冲击类型]
\normalsize
\textbf{① 气候冲击（产量）}：玉米、大豆的国内产量由播种面积与单产决定，单产受生长期
（东北 7---8 月）降水与积温影响；大豆、玉米的主产区高度重合，因此\emph{同一场天气会同时影响两个品种}，
但它们对价格的影响并不同步（大豆国内产量占比低，气候对价格的影响远弱于进口政策）。\par\vspace{3pt}
\textbf{② 政策冲击（进口与种植结构）}：玉米进口政策收紧、大豆扩种政策、玉米大豆生产者补贴、
三大主粮完全成本保险与种植收入保险全国覆盖\cite{cccp2024doc1}，
构成“国内价格地板”；而中央一号文件“深入推进粮油作物大面积单产提升行动”
与“种业振兴”\cite{cccp2025doc1}则决定长期的供给斜率。\par\vspace{3pt}
\textbf{③ 外生冲击（贸易与疫情）}：中美贸易摩擦对美豆份额的影响（2025 年美豆占比降至 15\%）、
以及非洲猪瘟这类疫情对\emph{下游需求}的冲击，是价格短期剧烈波动的主要来源。
\end{keybox}

\subsection{需求端：饲料工业的规模与结构}
\label{sec:grain-demand}

粮食与油料的需求侧可以用一个数字概括：\hl{2025 年全国工业饲料总产量 34225.3 万吨，
同比增长 8.6\%，创历史新高}\cite{chinafeed2025}。分品种看，
猪饲料 16639.4 万吨（+15.6\%，占总量的 48.6\%）、肉禽料 10097.7 万吨（+3.5\%，首次突破 1 亿吨）、
蛋禽料 3282.0 万吨（+1.4\%）、水产料 2323.1 万吨（+2.7\%）、反刍料 1475.8 万吨（+1.8\%）、
宠物料 188.4 万吨（+17.9\%）\cite{chinafeed2025}。上述六类合计 34006 万吨、占总量 \num{99.4}\%，
其余为其他饲料。

\begin{figure}[htbp]
\centering
\begin{tikzpicture}
\begin{axis}[agri axis, yearticks,
  width=13.6cm, height=5.8cm,
  xlabel={年份}, ylabel={全国工业饲料产量（万吨）},
  xmin=2022.6, xmax=2026.2, ymin=12000, ymax=35000,
  xtick={2023,2024,2025}, ytick={15000,20000,25000,30000,35000},
  % 图例移到图外上方：两条序列（总产量、猪饲料）数值区间不同，图内四角都会被数据占用
  legend style={at={(0.5,1.02)}, anchor=south, legend columns=2, cells={anchor=west}},
  legend entries={工业饲料总产量, 其中：猪饲料},
]
\addplot[AgriGreen, thick, mark=*, mark size=2pt] coordinates {
  (2023,32162.7) (2024,31503.1) (2025,34225.3)};
\addplot[HogPink, thick, mark=square*, mark size=2pt] coordinates {
  (2023,14975.2) (2024,14391.3) (2025,16639.4)};
\node[font=\scriptsize, InkGray, anchor=south] at (axis cs:2024,31700) {-2.1\%（近十年首降）};
\node[font=\scriptsize, AgriGreen, anchor=south] at (axis cs:2025,34350) {+8.6\%};
\end{axis}
\end{tikzpicture}
\caption{全国工业饲料产量与猪饲料产量（2023---2025）}
\label{fig:feed-output}
\end{figure}

饲料产量的波动本身就是一次“需求侧周期”的完整案例：2024 年受养殖业深度亏损影响，
全国饲料产量出现近十年首次下降（-2.1\%，其中猪料 -3.9\%、反刍料 -13.3\%、水产料 -3.5\%）\cite{chinafeed2024}；
2025 年在能繁母猪与家禽存栏高位的支撑下，猪料大幅回升 15.6\%，带动总量创新高\cite{chinafeed2025}。
需求侧的第二个结构性变化是\hl{豆粕减量替代}：2025 年饲料企业中豆粕在配合饲料和浓缩饲料中的
占比为 13.4\%、与上年持平，菜粕与棉籽粕等杂粕用量增加 3.0\%\cite{chinafeed2025}；
2023 年豆粕用量曾下降 11.8\%、2024 年再降 4.7\%\cite{chinafeed2024}。
这一趋势直接削弱了“生猪存栏$\to$豆粕需求$\to$豆粕价格”的传导强度。

\subsection{价格周期：玉米与豆粕的不同驱动}
\label{sec:grain-cycle}

\begin{figure}[htbp]
\centering
\begin{tikzpicture}
\begin{axis}[agri axis, yearticks,
  width=14.2cm, height=6.8cm,
  xlabel={年份（2021 年 2 月 $=100$）}, ylabel={归一化价格指数},
  xmin=2021.0, xmax=2026.9, ymin=35, ymax=200,
  xtick={2021,2022,2023,2024,2025,2026},
  % 图例移到图外下方：留在图内会遮住曲线末端的数值标注
  legend style={at={(0.5,-0.24)}, anchor=north, legend columns=5,
                cells={anchor=west}, font=\tiny},
  legend entries={生猪, 鸡蛋, 玉米, 豆粕, 豆油, 棕榈油, 白糖, 棉花, 天然橡胶},
]
\addplot[HogPink, very thick] table[x=x, y=LH]{data/clean/fut_norm_grain.dat};
\addplot[AgriGold, thick] table[x=x, y=JD]{data/clean/fut_norm_grain.dat};
\addplot[Grain, thick] table[x=x, y=C]{data/clean/fut_norm_grain.dat};
\addplot[OilBlue, thick] table[x=x, y=M]{data/clean/fut_norm_grain.dat};
\addplot[AgriGreen, thick] table[x=x, y=Y]{data/clean/fut_norm_grain.dat};
\addplot[violet, thick] table[x=x, y=P]{data/clean/fut_norm_grain.dat};
\addplot[SoftRed, thick] table[x=x, y=SR]{data/clean/fut_norm_grain.dat};
\addplot[AquaTeal, thick] table[x=x, y=CF]{data/clean/fut_norm_grain.dat};
\addplot[InkGray, thick] table[x=x, y=RU]{data/clean/fut_norm_grain.dat};
\node[font=\scriptsize, HogPink, anchor=east] at (axis cs:2026.88,196) {生猪 41.8};
\node[font=\scriptsize, Grain, anchor=east] at (axis cs:2026.88,180) {玉米 80.3};
\end{axis}
\end{tikzpicture}
\caption{主要农产品期货价格归一化走势（2021 年 2 月 $=100$，月度均价；图中数值为 2026 年 9 月的最新指数）}
\label{fig:fut-norm}
\end{figure}

图 \ref{fig:fut-norm} 呈现了本轮周期最关键的\hl{分化}：2021 年 2 月至 2026 年 9 月，
生猪期货价格下跌 58.3\%（指数 41.8），而玉米仅下跌 19.7\%、豆粕下跌 2.5\%、
豆油上涨 11.6\%、棕榈油上涨 43.8\%、天然橡胶上涨 24.8\%。
\emph{下游养殖品价格腰斩，而上游原料价格基本持平甚至上涨}——这直接说明
猪价下跌并未通过饲料需求传导为原料价格下跌，原因有三：一是生猪存栏的绝对水平并未崩塌
（2025 年生猪出栏 71973 万头、同比 +2.4\%），饲料需求仍在增长；
二是原料价格由国际供给与政策主导；三是豆粕减量替代降低了单位生猪的豆粕消耗。

\begin{figure}[htbp]
\centering
\begin{tikzpicture}
\begin{axis}[agri axis, catbar,
  width=13.0cm, height=7.4cm,
  xlabel={年化波动率（\%）}, xmin=0, xmax=30,
  yticklabels from table={data/clean/vol_bar.dat}{label},
  ytick=data,
  legend style={at={(0.97,0.03)}, anchor=south east},
  legend entries={主窗口 2021-02\textasciitilde2026-08, 全样本},
]
% 两组条形错开 0.2 行：同 y 绘制会互相遮盖，只剩较长的一根可见
\addplot[xbar, bar width=3.6pt, fill=AgriGreen!75, draw=AgriGreen]
  table[x=vol, y expr=\coordindex-0.2]{data/clean/vol_bar.dat};
\addplot[xbar, bar width=3.6pt, fill=AgriGold!80, draw=AgriGold]
  table[x=full, y expr=\coordindex+0.2]{data/clean/vol_bar.dat};
\end{axis}
\end{tikzpicture}
\caption{农产品期货年化波动率：主窗口与全样本对比（主窗口为 2021 年 2 月---2026 年 8 月的 67 个月度点；全样本自各品种上市首月至 2026 年 8 月；月度对数收益率年化，$\times\sqrt{12}$）}
\label{fig:vol-bar}
\end{figure}

波动率的分层同样值得注意（图 \ref{fig:vol-bar}）：生猪（27.1\%）、红枣（25.9\%）、
棕榈油（23.7\%）、鸡蛋（22.9\%）构成本轮周期的高波动组，
而粳米（4.0\%）、白糖（8.7\%）、玉米（9.2\%）波动率最低。
\hl{玉米的低波动是“政策定价”的直接结果}——进口受配额约束、国储与地方储备调节库存，
价格被限制在种植成本与替代品价格之间；而生猪、鸡蛋、红枣完全由市场供需定价，
且供给调整存在刚性，因此其波动率约为玉米的三倍（生猪 27.1\% / 玉米 9.2\%）。

\subsection{与猪周期的连接：量与价的分离}
\label{sec:grain-hog-link}

把上述供给与需求拼起来，可以得到本章的核心判断：

\begin{keybox}[粮食油料与猪周期的连接方式]
\normalsize
\textbf{（1）需求侧同向、价格侧弱相关。}饲料产量（尤其是猪料）与生猪存栏高度同向，
2025 年猪料增长 15.6\% 就是生猪存栏高位的直接结果；但玉米、豆粕的\emph{价格}
与生猪价格的相关性很弱（第 \ref{sec:corr} 章的月度收益率相关系数分别为
玉米 $+0.09$、豆粕 $+0.13$，均未达到 $|\rho|>0.24$ 的 5\% 显著性阈值），
因为原料价格由国际供给与政策决定。\par\vspace{3pt}
\textbf{（2）成本传导是非对称的。}原料价格上涨会立即压缩养殖利润（玉米与豆粕占饲料成本
超过 70\%）\cite{sun2026hog}，但养殖亏损压缩产能后，原料需求的下降幅度有限
（存栏的下行慢于利润的下行）。这构成养殖业“利润周期”比“价格周期”更剧烈的原因。\par\vspace{3pt}
\textbf{（3）政策是两条链的共同外生变量。}同一份中央一号文件既规定粮食产能目标
（“确保粮食产量保持在 1.3 万亿斤以上”）\cite{cccp2024doc1}，又规定生猪产能调控机制，
使粮食与养殖在\emph{数量}上被同时锁定、在\emph{价格}上被各自隔离。
\end{keybox}
